\chapter*{Introduction générale}
\addcontentsline{toc}{chapter}{Introduction générale}

La transition vers la mobilité électrique entraîne le déploiement d'un nombre croissant
de bornes de recharge. L'exploitation de ces équipements ne repose pas uniquement sur
leur installation physique : elle exige également une supervision continue, une
connaissance fiable de leur disponibilité et une coordination efficace entre les
utilisateurs, les opérateurs et les techniciens.

Dans ce contexte, le présent stage d'été, réalisé au sein de \CompanyName{}, porte sur
la conception et le développement de \ProjectName{}, une plateforme web consacrée à la
visualisation et au suivi de la disponibilité des bornes de recharge électrique. Le projet
associe gestion multi-organisation, échanges OCPP, cartographie, sessions de recharge,
paiement simulé, alertes, interventions et reporting.

Le rapport conserve l'organisation en cinq chapitres du document de référence. Le premier
chapitre introduit le cadre du projet, le domaine métier, la problématique, l'existant et
la démarche retenue. Le deuxième formalise les acteurs, les besoins et la planification.
Le troisième présente l'architecture et la conception détaillée. Le quatrième décrit la
réalisation incrémentale, l'intégration et la validation. Enfin, le cinquième illustre les
principales interfaces et les parcours adaptés aux différents rôles.

\plannedcontent{La version rédigée précisera les objectifs pédagogiques et professionnels
du stage, le périmètre validé avec l'encadrant, les contributions réalisées et les limites
du travail présenté.}

